%% conclusion.tex

\section{Conclusion}
\label{sec:conclusion}

\paragraph{Summary.}
This paper introduced the \emph{topology-mutation privilege-escalation} attack vector, a previously overlooked class of authorization failures on the recovery path of self-healing LLM agent systems. We showed that the recovery pipeline---designed to improve robustness by allowing a diagnostic LLM to propose remedial actions---can be abused as a \emph{privilege-escalation trampoline}: a recovery strategy that proposes to mutate the agent topology (e.g., to replace a role or reconfigure a sub-agent) bypasses the authorization checks that gate the forward execution path, because the recovery machinery was never designed to re-verify privileges at the side-effect boundary. The contribution is architectural: we formalize the anti-pattern as a manifestation of CWE-862 (Missing Authorization) on the recovery sink, trace its root cause to the classical \emph{confused deputy} problem, and propose a single OS-principled design pattern---the \reauthgate{}---that ports the kernel exception-handler privilege invariant to the LLM agent recovery path.

\paragraph{Architectural contributions.}
The paper makes three architectural contributions. First, we identify and formalize the \emph{unchecked recovery sink} anti-pattern, in which a recovery strategy executes a privileged side effect (role replacement, topology mutation) without re-checking the authorization that gates the forward path; we trace its root cause to the implicit ``failure = benign, retry is safe'' assumption inherited from classical OS exception handling---an assumption that holds when failure payloads originate from the trusted computing base but breaks down when a diagnostic LLM proposes privileged remedial actions. Second, we propose the \reauthgate{}, a single design pattern that enforces the \privinv{} ($P(\text{recovery action}) \leq P(\text{failed execution})$) at a centralized choke point through which all side-effecting recovery actions must flow; the pattern ports the OS principle that exception handlers never escalate privileges and the capability-based security tradition to the agent recovery path. Third, we show that the pattern is framework-agnostic: it depends only on the existence of a recovery pipeline with declarable side effects and a permission checker, not on any specific agent architecture, making it portable to any self-healing agent system that exposes a topology-mutation or role-replacement primitive.

\paragraph{Validation summary.}
A concise validation on the \system{} testbed confirms that the \reauthgate{} is implementable, effective, and imposes negligible cost on legitimate self-healing. The defended V2 attack success rate is \textbf{0.00 by construction}: the whitelist invariant rejects out-of-set roles deterministically, illustrating the key architectural principle that structural invariants enforced at a centralized choke point are LLM-independent and therefore robust to model scaling. The defense imposes zero false alarms on legitimate recovery and topology-mutation workloads, and sub-millisecond machinery overhead. Because the guarantee is structural rather than probabilistic, it does not degrade as instruction-following models become stronger, and it is not bypassable by paraphrased payloads or defense-aware adversaries---the gate inspects the proposed role against an authorization set, not the language of the diagnostic recommendation.

\paragraph{Empirical grounding.}
The empirical evidence in this paper comprises a source-level structural analysis of the \system{} recovery pipeline that identifies the \code{parseAndValidate(validate=false)} call in \code{TopologyMutationStrategy} as a CWE-862 missing-authorization site, and a measured end-to-end attack that achieves ASR 1.00 in the baseline configuration and ASR 0.00 in the defended configuration across $N=30$ trials. The structural argument establishes that the recovery-path authorization boundary is independent of forward-path input defenses: even with forward-path input sanitizers, instruction hierarchies, or capability-based tool invocation fully enabled, the attack would still succeed because the privileged side effect is invoked by the recovery orchestrator, which forward-path defenses never inspect. This confirms that topology-mutation privilege escalation is a distinct authorization surface on the recovery path, not a variant of indirect prompt injection.

\paragraph{Limitations and future work.}
The \reauthgate{} enforces a \emph{static} privilege invariant: it checks that the proposed role is in the legitimate downgrade set, but does not inspect \emph{why} the recovery strategy proposed it. A sufficiently sophisticated payload could induce a strategy to propose an action that is individually permissible yet collectively harmful (e.g., a sequence of individually-legitimate downgrades that compose into an unintended capability). As future work, we propose \emph{intent auditing via inner-monologue inspection}: instrumenting the diagnostic LLM's chain-of-thought to verify that the proposed recovery action is justified by the failure's actual cause, rather than by injected reasoning. This would add a semantic, model-level defense layer atop the syntactic, OS-level invariant introduced here---complementary rather than a replacement, and aligned with the broader trend of using LLMs to audit LLM behavior.

A second limitation concerns the scope of the validation: the testbed validation establishes implementability, efficacy, and cost, but broader empirical claims---cross-model robustness of the content-level diagnostics that feed the gate, additional production-framework topology-mutation primitives, and fully adaptive multi-round adversaries that observe defense rejections and iteratively rewrite diagnostic recommendations---are not established in this architectural paper and are immediate future work.
